\documentclass[11pt,a4paper,addpoints]{exam}
\usepackage[utf8]{inputenc}
\usepackage[T1]{fontenc}
\usepackage[polish]{babel}
\usepackage{geometry}

\pointpoints{pkt}{pkt}
\boxedpoints
\hqword{Zadanie:}
\hpword{Punkty:}
\hsword{Wynik:}
\htword{Suma}


\lhead{\makebox[12cm][l]{\large\textbf{Imię i nazwisko:}\enspace\hrulefill}}
\rhead{W02: msbuild}

\footer{}{}{}

\begin{document}

\cellwidth{2.75cm}
\noindent \gradetable[h]

\vspace{0.8cm}

Zdecydowałeś się stworzyć aplikację służącą do wykonywania prostych operacji związanych z ciągiem Fibonacciego.
Aplikacja powinna umożliwiać wyznaczanie kolejnych wyrazów tego ciągu.

\begin{questions}

\question[5]
Utwórz nową solucję zawierające następujące trzy projekty:

\begin{itemize}
    \item \texttt{Fibonacci.Lib} -- projekt biblioteki,
    \item \texttt{Fibonacci.App} -- aplikację konsolową,
    \item \texttt{Fibonacci.Tests} -- projekt z testami.
\end{itemize}

Dodaj odpowiednie referencje pomiędzy projektami.

W projekcie biblioteki utwórz klasę zawierającą metodę:
\newline
\texttt{GetFibonacci}, która zwraca wyraz ciągu Fibonacciego o podanym numerze,


Metoda \texttt{GetFibonacci} powinna przyjmować jako argument liczbę całkowitą określającą numer wyrazu ciągu i zwracać liczbę całkowitą.

Przyjmij, że:

\[
F_0 = 0,\qquad F_1 = 1,\qquad
F_n = F_{n-1} + F_{n-2}.
\]


W aplikacji konsolowej utwórz pętlę \texttt{for}, która wypisze 10 pierwszych wyrazów ciągu Fibonacciego.
Do obliczenia każdego wyrazu wykorzystaj metodę \texttt{GetFibonacci} z projektu biblioteki.

W projekcie testowym utwórz trzy testy tej funkcji.

Testy powinny sprawdzić, czy \texttt{GetFibonacci} zwraca poprawne wartości dla wybranych numerów wyrazów ciągu,


\end{questions}

\end{document}